\section*{Chapter \ref{ch:dv:fx-derivatives} --- FX Derivatives}

\begin{solution}{exo:dv:fx-derivatives:1}
$e^{-r_fT}S=e^{-r_dT}F$, so $e^{-r_fT}S\Phi(d_1)-e^{-r_dT}K\Phi(d_2)=e^{-r_dT}\bigl(F\Phi(d_1)-K\Phi(d_2)\bigr)$, Black's formula on the forward with
the quote currency's discount factor; $d_1$ and $d_2$ depend on $S$ only through $F$.
\end{solution}

\begin{solution}{exo:dv:fx-derivatives:2}
The market's fear is a fall of the emerging currency, a rise of the dollar against it, and such falls come with rising volatility: the rate and its
variance move together, $\rho>0$. Calls on the dollar (the protection) are then priced at higher volatilities than puts, and the risk reversal is
positive: $+1.89$ points at 25 delta.
\end{solution}

\begin{solution}{exo:dv:fx-derivatives:3}
A leveraged short option position: on each fixing above the strike the client sells twice the notional at a rate below the market, and the target caps
the gains; the improvement is the premium of that position, received as a better rate.
\end{solution}

\begin{solution}{exo:dv:fx-derivatives:4}
The weights match the pillars' \omterm{def:dv:greeks-and-the-hedging-pnl:greeks}{vega}, \omterm{def:dv:greeks-and-the-hedging-pnl:greeks}{vanna} and \omterm{def:dv:greeks-and-the-hedging-pnl:greeks}{volga} exactly, so at a pillar the weight is one on itself and zero elsewhere, and the price is the market
price. Beyond the pillars the method assumes the smile's cost is carried by those three risks at the at-the-money volatility; the 10-delta wings have a
curvature the 25-delta quotes do not pin down, so the extrapolation misses it (a point at the 10-delta call here).
\end{solution}

\begin{solution}{exo:dv:fx-derivatives:5}
In \omterm{def:dv:local-volatility:model}{local volatility} the volatility is a function of the spot, and the market's call wing is expensive, so as the rate rises towards 7.20 its volatility
rises with it and it reaches the barrier more easily. In Heston part of the volatility is carried by a variance process that is not tied to the spot, so
the path to the barrier is less volatile on average.
\end{solution}

\begin{solution}{exo:dv:fx-derivatives:6}
With $\lambda=0$ the variance follows $dv=\kappa(\bar v-v)dt$, a deterministic path $v(t)=\bar v+(v_0-\bar v)e^{-\kappa t}$; then
$\E[v_t\mid S_t=S]=v(t)$ and $L(t,S)^2v(t)=\sigma_{\mathrm{loc}}(t,S)^2$.
\end{solution}

\begin{solution}{exo:dv:fx-derivatives:7}
The zero-cost strike is 7.325 under \omterm{def:dv:local-volatility:model}{local volatility}, 7.312 under stochastic-local volatility and 7.305 under Heston; the expected number of fixings is 1.64,
1.73 and 1.78. The model with the more volatile path to the upper side (\omterm{def:dv:local-volatility:model}{local volatility}) prices the client's leveraged short calls higher, so the dealer
can offer a better strike.
\end{solution}

\begin{solution}{exo:dv:fx-derivatives:8}
The 98\% redeem with a small, capped gain (0.30 per dollar at most); the other 2\% run with twice the notional against the client and carry all the product's
value: the expected loss in those paths is large enough to make the whole worth zero. Probability of redemption says nothing about the size of the losses.
\end{solution}

\begin{solution}{pb:dv:fx-derivatives:1}
\textbf{1.} 5.32\%; $+1.89$ and $+0.46$ points.
\textbf{2.} That the dollar's rise against the currency, a devaluation, is the feared scenario, priced at higher volatility.
\textbf{3.} 6.861: the dollar's rate (5\%) is above the currency's (3\%), so the forward is below the spot by the carry.
\textbf{4.} Within a tenth of a point at 10 and 25 delta and at the money.
\textbf{5.} 0.79.
\textbf{6.} 7.312, 6.6\% above the forward.
\textbf{7.} 98.1\% redeem, after 1.73 fixings on average.
\textbf{8.} A gain of 0.30 million within one or two fixings on most paths (95\% gain at least 126\,000); the zero value comes from the few paths on which the
rate rises and stays above the strike, costing twice the notional every month.
\textbf{9.} 7.325 under \omterm{def:dv:local-volatility:model}{local volatility}, 7.305 under Heston.
\textbf{10.} A larger target lets the client earn more before redemption, so the product lives longer and exposes the dealer less on the leveraged side for a
given strike; to stay at zero cost the strike must be less generous.
\textbf{11.} 7.12 million per million of monthly notional, of which 1.01 million in the first quarter.
\textbf{12.} 0.99 million.
\textbf{13.} 11.7 fixings: after a 10\% rise it almost never redeems.
\textbf{14.} Its delta changes sharply: close to the target the remaining gains are small and the product is about to vanish, so the hedge of the future
leveraged forwards must be unwound; the delta can flip sign within a few pips.
\textbf{15.} A step depreciation of the emerging currency in a managed regime is the scenario on which the TARF's value rests: a 2.8\% fall in two days after a
change of fixing rule is of the kind that turns a product expected to redeem at once into one that runs its course at twice the notional.
\textbf{16.} The distribution of outcomes, not just the redemption probability: the loss in the scenarios that matter (a 10\% move costs 7.1 million per
million), and the leverage.
\textbf{17.} The dynamics of the smile and jumps: the zero-cost strike moves by 0.02 (200 pips) between models, but a step devaluation is outside all three.
\textbf{18.} No: without a natural dollar income the leveraged sales are speculative, not a hedge; a regulator and the bank's suitability rules would object.
\textbf{19.} 1.73 fixings on average before redemption; an expected loss of 7.12 million per million of monthly notional if the rate rises 10\% over the first
quarter.
\textbf{20.} A leveraged strip of options on the rate moving against it, paid for with a better rate and capped gains.
\end{solution}

\begin{solution}{iq:dv:fx-derivatives:1}
At the money: the delta-neutral straddle's volatility. Risk reversal: the 25-delta call's volatility minus the put's (the skew). Butterfly: the wings'
average over at the money (the curvature), in the pair's convention (smile or broker strangle). Convert with the pair's delta convention into three
strike--volatility points and interpolate (vanna--volga, SABR by expiry, or a parametric smile).
\iqlookfor{the three quotes and the conventions.}
\end{solution}

\begin{solution}{iq:dv:fx-derivatives:2}
Solve for the weights of the three pillars matching the option's \omterm{def:dv:greeks-and-the-hedging-pnl:greeks}{vega}, \omterm{def:dv:greeks-and-the-hedging-pnl:greeks}{vanna} and \omterm{def:dv:greeks-and-the-hedging-pnl:greeks}{volga} at the at-the-money volatility; add the weighted market-minus-Black
costs of the pillars to the Black price. It assumes the smile's cost is fully captured by those three risks priced by the quoted instruments, a ``flat but
stochastic'' volatility.
\iqlookfor{the weights and the assumption.}
\end{solution}

\begin{solution}{iq:dv:fx-derivatives:3}
$dS/S=\dots+L(t,S)\sqrt v\,dW$ with a stochastic variance: the \omterm{def:dv:fx-derivatives:leverage}{leverage function} makes it fit every vanilla, the variance process supplies dynamics.
Calibrate $L(t,S)^2=\sigma_{\mathrm{loc}}^2/\E[v\mid S]$ by the particle method: simulate, estimate the conditional expectation from the particles at each step,
update $L$; choose the mixing of the volatility of variance from exotic prices.
\iqlookfor{the projection formula and the particle method.}
\end{solution}

\begin{solution}{iq:dv:fx-derivatives:4}
The difference is in the dynamics, not the smile: calibrate the model's dynamic parameter (the mixing weight) to the market's touch prices where they trade;
otherwise price with the model whose dynamics the historical behaviour of the smile supports, and hold a reserve for the spread between the plausible
models.
\iqlookfor{calibrate dynamics to exotics, and reserve.}
\end{solution}

\begin{solution}{iq:dv:fx-derivatives:5}
The client sells a currency forward at a better rate while he wins, with gains capped by a target, and twice the amount at a worse rate while he loses, with
no cap. For the client: large losses in a sustained move against him. For the bank: the hedge's sensitivity to the smile's dynamics and to jumps, and the
client's credit and suitability risk once the losses mount.
\iqlookfor{the asymmetry, and both sides' risks.}
\end{solution}

\begin{solution}{iq:dv:fx-derivatives:6}
Add a jump (or regime-change) component to pricing and to the scenarios: barriers can be crossed without trading near them, and a TARF can go from redeeming
to running at double notional overnight. Hedge with options that pay in the jump, limit concentration at barrier levels near a band's edge, and reserve for the
jump at the book level.
\iqlookfor{jumps in pricing, gap hedges and reserves.}
\end{solution}
